% Методы (main.tex)
\section{Методы}\label{sec:methods}

\subsection{Обозначения и постановка}

Блок Qwen3~\cite{vaswani2017attention} с предварительной нормализацией
действует на остаточный поток $h\in\mathbb{R}^{T\times d}$ как
\[
  h' = h + \mathrm{Attn}(\hat x),\quad \hat x=\mathrm{RMSNorm}(h),\qquad
  h'' = h' + \mathrm{MLP}(\hat x'),\quad \hat x'=\mathrm{RMSNorm}(h').
\]
Внимание групповое (GQA): $H$ голов запросов делят $H_{kv}$ голов
ключей и значений, размерность головы $D$. Запросы и ключи проходят
нормализацию по размерности головы (QK-norm) и поворотное
позиционное кодирование (RoPE):
$q_t=\mathrm{RoPE}_t(\mathrm{RMSNorm}_q(W_q\hat x_t))$,
$k_j=\mathrm{RoPE}_j(\mathrm{RMSNorm}_k(W_k\hat x_j))$, $v_j=W_v\hat x_j$,
$a_{tj}=\mathrm{softmax}_{j\le t}(q_t^\top k_j/\sqrt D)$.
MLP имеет вид SwiGLU: $\mathrm{MLP}(x)=W_d\,z(x)$,
$z(x)=\mathrm{silu}(W_g x)\odot W_u x\in\mathbb{R}^{I}$.

Заменяются слои $0,\dots,k-1$. Различаем три вида замены: только
подслой внимания (MLP точный), только MLP (внимание точное) и блок
целиком. Детерминированный метод строится по весам и по
калибровочному набору $\mathcal C$ (фрагменты WikiText-103) в замкнутой
форме: средние, ковариации, SVD, гребневая регрессия и
поочерёдная минимизация (backfitting). Стохастическая аппроксимация
и градиентное обучение в основном треке не используются.

\paragraph{Позиция 0.} В каузальной модели состояние первого токена
зависит только от самого токена, а любая нормированная строка внимания
при $t=0$ даёт ровно $W_o v_0$. Поэтому замены внимания на позиции $0$
точны автоматически. Для замен MLP и блоков мы вычисляем позицию $0$
исходным слоем, что эквивалентно таблице по словарю из $|V|$ строк.
Это существенно: в Qwen3 норма состояния первого токена вырастает на
два–три порядка после нескольких первых слоёв (массивная активация,
на которой держится сток внимания), и любая ошибка в ней разрушает
модель.

\subsection{Методы из литературы, адаптированные к Qwen3}

\paragraph{Внимание.}
\textbf{L1} (окно + сток; Sparse Transformer~\cite{child2019sparse},
BigBird без случайных рёбер~\cite{zaheer2020bigbird}):
softmax вычисляется только по $j$ с $t-j<w$ или $j=0$.
\textbf{L2} (strided~\cite{child2019sparse}): дополнительно
разрешены $j$ с $(t-j)\bmod s=0$.
\textbf{L3} (Performer, FAVOR+~\cite{choromanski2021performer}):
$\exp(q^\top k/\sqrt D)\approx\varphi(q')^\top\varphi(k')$, $q'=qD^{-1/4}$,
$\varphi(x)=m^{-1/2}\exp(\Omega x-\|x\|^2/2)$ с блочно-ортогональной
гауссовой $\Omega\in\mathbb{R}^{m\times D}$; это единственный
стохастический метод сравнения, он запускается с тремя seed.
\textbf{L4} (линейное внимание~\cite{katharopoulos2020rnn}):
$\varphi(x)=\mathrm{elu}(x)+1$.
\textbf{L5} (фиксированное внимание в духе
Synthesizer~\cite{tay2021synthesizer}): $a_{tj}$ зависит только от
$t-j$: ядро $K_h(\delta)$ при $\delta<\Delta$, вес стока $\sigma_h$ при
$j=0$, равномерный <<хвост>>; $K_h$ и $\sigma_h$ "--- средние точного
внимания на $\mathcal C$, строки перенормируются.
\textbf{L6} (A3~\cite{wong2025a3}). Код авторов поддерживает Llama и
MPT, поэтому мы перенесли метод на Qwen3. QK: в каждой KV-группе
оставляются $r/2$ пар RoPE $(i,i+D/2)$ с наибольшим
$\sum_{h\in g}\sqrt{\mathbb E q_{h,i}^2}\sqrt{\mathbb E k_{g,i}^2}$; моменты
берутся после QK-norm, что совпадает с критерием A3
$\|R^{1/2}w_q\|\,\|w_k R^{1/2}\|$ для линейных проекций. VO: для группы
$g$ строится $M_g=[W_{v,g}^\top W_{o,h}^\top]_{h\in g}$ и усечённое
SVD $R^{1/2}M_g=U\Sigma V^\top$, где $R=\mathbb E\,\hat x\hat x^\top$. Мы
полагаем $W_v'^\top=M_gV_r$, $W_o'=V_r^\top$. Произведение $M_gV_rV_r^\top$
совпадает с решением A3 при обратимой $R$, но не требует $R^{-1/2}$,
которая на реальных данных плохо обусловлена. FFN: оставляются нейроны
с наибольшим $\sqrt{\mathbb E z_i^2}\,\|w_{d,i}\|$.

\paragraph{MLP.}
\textbf{ML1} "--- отбор нейронов A3 (см. выше).
\textbf{ML2} "--- <<оракульная>> контекстная разреженность Deja
Vu~\cite{liu2023dejavu}: на каждом токене используются $\kappa$
нейронов с наибольшим $|z_i|$. Это верхняя граница для обучаемого
предиктора Deja Vu и одновременно разреженность памяти
<<ключ--значение>>~\cite{geva2021ffn}.
\textbf{ML4} "--- ASVD~\cite{yuan2023asvd}: $S=\mathrm{diag}(\mathbb E|x|)^{\alpha}$,
$W\approx(U_r\Sigma_r)(V_r^\top S^{-1})$ по SVD матрицы $WS$, $\alpha=0.5$.
\textbf{ML5} "--- SVD-LLM~\cite{wang2025svdllm} без шага дообучения:
$\mathbb E\,xx^\top=LL^\top$ (Холецкий), $W\approx U_r\Sigma_rV_r^\top L^{-1}$
по SVD матрицы $WL$. Для $W\in\mathbb R^{m\times n}$ и доли параметров
$\rho$ ранг $r=\lfloor\rho mn/(m+n)\rfloor$; сжимаются $W_g$, $W_u$, $W_d$.
\textbf{ML7} "--- MoEfication~\cite{zhang2022moefication}: нейроны
разбиваются сбалансированным $k$-means (детерминированная инициализация
$k$-means++) по нормированным строкам $[W_g;W_u]$ на $E$ экспертов;
маршрутизация без обучения выбирает $m$ экспертов с наибольшим
$\hat x^\top c_e$, где $c_e$ "--- центроид строк $W_g$ эксперта.

\paragraph{Неприменимые методы.}
Linformer~\cite{wang2020linformer} проецирует ключи и значения по оси
последовательности обучаемыми матрицами. Nyströmformer~\cite{xiong2021nystrom}
строит опорные точки как средние по сегментам всей
последовательности, FNet~\cite{leethorp2022fnet} применяет ДПФ ко всем
позициям. В каузальной модели все три метода смешивают будущие токены
с прошлыми, и маскирование после факта не восстанавливает
каузальность. SOFT~\cite{lu2021soft} заменяет softmax гауссовым ядром,
Reformer~\cite{kitaev2020reformer} требует общей проекции QK и
случайного хеширования; оба предполагают обучение с нуля.

\subsection{Предложенные методы}

\paragraph{O1: типизированная замена голов.}
Для каждой головы вычисляется расстояние полной вариации
$\mathrm{TV}_h=\mathbb E_t\tfrac12\|a_{t\cdot}-\tilde a_{t\cdot}\|_1$ между
точным вниманием и позиционным ядром L5 (на калибровке). Доля $\rho$
голов с наименьшим $\mathrm{TV}_h$ заменяется ядром, остальные считаются
точно (или в окне $w$). Ядро не требует вычисления $q^\top k$.

\paragraph{O2: ридж-свёртка вместо внимания.}
Подслой внимания заменяется каузальной линейной свёрткой в пространстве
$p$ главных компонент $P$ входа:
$\mathrm{Attn}(\hat x)_t\approx\sum_{\delta<w}W_\delta P^\top\hat x_{t-\delta}+W_sP^\top\hat x_0+b$.
Коэффициенты находятся гребневой регрессией.

\paragraph{O3: аддитивная $n$-граммная замена блоков.}
Выход первых $k$ блоков при $t\ge1$ приближается суммой таблиц по
токенам окна и позиционного члена:
$h_{k,t}\approx e(x_t)+\sum_{\delta<m}T_\delta[x_{t-\delta}]+\pi[\min(t,P)]$.
Таблицы подгоняются поочерёдной минимизацией (backfitting) со сжатием
$T[v]=\sum r/(n_v+\lambda)$. Вычисление первых $k$ слоёв заменяется $m$
обращениями к таблице.

\paragraph{O4: позиционное ядро $\times$ значимость ключа.}
Профиль (раздел~\ref{sec:results}) показывает, что внимание первых слоёв
\emph{контентное}: позиционное ядро объясняет его хуже, чем в средних
слоях. Мы добавляем к ядру множитель, зависящий от токена-ключа:
\[
  a_{tj}\;\propto\;K_h(t-j)\,\exp\bigl(\beta_h(x_j)\bigr),\quad j\ge1;\qquad
  a_{t0}\propto\sigma_h .
\]
Таблица $\beta_h$ подгоняется итеративным согласованием маргиналов
(IPF): $\beta_h(v)\leftarrow\beta_h(v)+\log\frac{A^{\mathrm{true}}_h(v)+\alpha}{A^{\mathrm{model}}_h(v)+\alpha}$,
где $A_h(v)$ "--- суммарная масса внимания, полученная ключами с токеном
$v$ на $\mathcal C$. Как и L5, метод не вычисляет $q^\top k$, но
улавливает головы, смотрящие на определённые токены (пунктуацию,
служебные слова). В гибридном варианте доля голов с наибольшей
остаточной TV-ошибкой считается точно (при необходимости в окне $w$).

\paragraph{M-O1: линеаризация MLP с точными нейронами.}
$\mathrm{MLP}(\hat x)\approx W\hat x+b+\sum_{i\in S}w_{d,i}z_i(\hat x)$,
где $S$ "--- $r$ нейронов по критерию A3, а $W,b$ "--- гребневая
регрессия на остаток. При $S=\varnothing$ MLP заменяется одной
линейной картой ($d^2$ параметров вместо $3dI$).

\paragraph{M-O4: предиктор активных нейронов в замкнутой форме.}
Контекстная разреженность Deja Vu~\cite{liu2023dejavu} использует
обучаемый предиктор активных нейронов. Мы заменяем его гребневой
регрессией: вклад нейрона $|z_i(\hat x)|\,\|w_{d,i}\|$ предсказывается
линейно из $p=64$ главных компонент входа, $\widehat{c}(\hat x)=(\hat xP)A$.
На каждом токене вычисляются только $\kappa$ нейронов с наибольшим
$\widehat c_i$, а остаток MLP компенсируется линейной картой $W\hat x+b$,
подогнанной на остаток. Стоимость "--- $3d\kappa+d^2+p(d+I)$ умножений
вместо $3dI$. Вариант с таблицей $N[x_t]$ (набор нейронов только по
текущему токену) заметно хуже: активность нейронов первых слоёв
зависит от контекста.

\paragraph{O5: окно + детерминированный хвост.}
Внутри окна $w$ и на стоке оценки $q^\top k$ точные, для дальних ключей
оценка заменяется таблицей $\beta_h(x_j)$, подогнанной IPF по массам
внимания. Хвост вычисляется префиксными суммами за $O(T)$.
Инициализация $\beta$ лог-средним $\exp(s)$ завышает массу хвоста
(неравенство Йенсена) и ухудшает качество в разы; со средним $s$ и IPF
метод сравним с простым окном того же размера. Это указывает, что
дальний контекст в первых слоях почти не используется.

\paragraph{M-O2 и M-O3.} Табличный MLP:
$\mathrm{MLP}(\hat x_t)\approx T[x_t]+W\hat x_t+b$ (поочерёдная подгонка
таблицы и линейной части). Кусочно-линейный MLP: $k$-means в
пространстве главных компонент входа и своя гребневая карта в каждом
кластере.

\paragraph{Последовательная калибровка с компенсацией дрейфа.}
Все наши методы подгоняются слой за слоем. Признаки слоя $i$
вычисляются по пути <<ученика>>, в котором слои $<i$ уже заменены, а
целью служит состояние учителя: $y=h^{T}_{i+1}-h^{S}_{i,\mathrm{mid}}$.
Тем самым каждый заменённый слой не только аппроксимирует свой
исходный, но и возвращает остаточный поток к траектории исходной
модели. Методы из литературы калибруются один раз по пути исходной
модели, как в оригинальных работах.

\subsection{Дистилляционный контроль}
Чтобы сравнить с методами, которые в оригинале обучаются (линейное
внимание, шаг дообучения SVD-LLM), и оценить запас наших методов, мы
отдельно дообучаем параметры замен AdamW (300 шагов) на потерю
$\mathbb E_t\|h^S_{k,t}-h^T_{k,t}\|^2/\|h^T_{k,t}\|^2$. Эти результаты не
детерминированы в смысле введения и в таблицах помечены отдельно.
